\documentclass[11pt]{letter}
\usepackage[margin=0.9in]{geometry}
\usepackage{hyperref}
\usepackage{enumitem}
\usepackage{parskip}

% Tighten letter-class vertical spacing to fit one page.
\setlength{\parskip}{4pt}
\setlength{\longindentation}{0pt}

\signature{Guilin Zhang, Kai Zhao, Xu Chu}
\address{}
\date{}

\begin{document}

\begin{letter}{The Editors-in-Chief\\ \textit{Quantum Machine Intelligence}\\ Springer Nature}

\opening{Dear Editors,}

We are pleased to submit our manuscript, ``Budget-Aware Neural Error Mitigation for Variational Quantum Algorithms,'' by Guilin Zhang, Kai Zhao, and Xu Chu, for consideration as a Research Article in \textit{Quantum Machine Intelligence}. The journal's Aims \& Scope explicitly lists ``AI for Quantum Error Correction and Mitigation'' as a topic of interest, and our work sits squarely within this remit: we use a learned neural model to mitigate noise in variational quantum algorithms (VQAs).

The work introduces a budget-conditioned neural error mitigation (NEM) model that corrects a single shot-limited noisy VQA expectation value from one quantum circuit execution, conditioned on circuit features, device-characterization features, and the shot-noise scale $1/\sqrt{S}$, so that one network serves all shot budgets. Crucially, we evaluate it under an equal-budget protocol in which every method---NEM, zero-noise extrapolation (ZNE), and Clifford data regression (CDR)---spends the same total number of shots per evaluation point, split across the executions it requires (NEM~1, ZNE~3--5, CDR~32). This is the comparison that VQA optimization loops actually face, and it differs from prior neural-QEM work, which grants every method full per-execution shots.

Our three contributions are:
\begin{itemize}[leftmargin=*,topsep=2pt,itemsep=2pt]
  \item \textbf{Budget-conditioned neural mitigation:} one network, one execution per evaluation, serving all shot budgets via $1/\sqrt{S}$ conditioning.
  \item \textbf{An equal-budget evaluation protocol and Pareto analysis} that reflects real deployment cost---the methodological core of the paper.
  \item \textbf{A characterization of when extrapolation-based ZNE catastrophically fails}---returning a worse answer than no mitigation on 38--65\% of miscalibration instances, with unbounded single-instance error---under physically motivated coherent miscalibration.
\end{itemize}

Under combined stochastic and coherent control-offset miscalibration noise, NEM is the only method that beats unmitigated execution at budgets up to $2^{14}$ shots ($+72$--$81\%$); CDR overtakes only at $2^{16}$ shots, at $32\times$ the per-evaluation quantum cost. These improvements hold from 4 to 20 qubits. We are deliberate about what we do not claim: CDR is more accurate wherever its $32\times$ per-evaluation cost is affordable, and our architecture is a plain residual MLP---the contribution is the protocol, the budget-conditioning, and the failure analysis, not a new network. We regard this candor as a feature of the work's scientific integrity.

All simulation code and per-instance result data are public at \url{https://github.com/GuilinDev/Quantum-Error-Mitigation}, and the manuscript includes a complete Declarations section with an AI-use statement.

% TODO: include once hardware run is done -- delete this entire block if hardware is not added.
%To strengthen the empirical claims, we additionally include a small real-device validation on IBM quantum hardware that reproduces the central ordering---ZNE fails while NEM holds---under realistic coherent miscalibration, confirming that our simulation-based conclusions transfer to physical devices.

We confirm that this manuscript is not under consideration elsewhere and that the authors have no competing interests to declare. We believe the work is a strong fit for the journal's scope and would be grateful for your consideration. Thank you for your time.

\closing{Sincerely,}

\end{letter}

\end{document}
